\section{Preliminaries}

\subsection{BDI Agents and Plan Libraries}

[inline def] Represent a BDI achievement plan as
\[
r \;=\; +!p(\bar X):C \leftarrow b_1;\ldots;b_m,
\]
where \(+!p(\bar X)\) is the achievement-goal event handled by the plan,
\(C\) is its applicability context, and
\(\beta_r=\langle b_1,\ldots,b_m\rangle\) is its plan body.

[inline def] Let each body step \(b_i\) be either a primitive domain action
or a subsidiary achievement-goal call \( !q(\bar Y) \).

[inline def] Define a BDI plan library
\[
\mathcal L=\{r_1,\ldots,r_k\}
\]
as a finite set of BDI plans available to the agent.

[inline def] Let a grounding substitution \(\theta\) map plan variables to
type-compatible domain objects, and write \(r\theta\) for the corresponding
ground instance of plan \(r\).

[inline def] Call \(r\) parameterised when it contains variables that can be
grounded to different domain objects, rather than being tied to one particular
planning instance.

[inline def] Define a plan \(r\theta\) as applicable in state \(s\) when its
grounded context holds, written
\[
s \models C\theta.
\]

[inline def] Define a subsidiary call \( !q(\bar Y) \) as invoking an
applicable plan in \(\mathcal L\) whose triggering achievement goal
type-compatibly matches \(q(\bar Y)\).


\subsection{STRIPS Planning and Generalized Planning}

[formal def] Define the supported planning domain as a typed STRIPS domain
with bounded-integer resources,
\[
D=\langle \mathcal P,\mathcal F,\mathcal A,\mathcal T\rangle,
\]
where \(\mathcal P\) is a finite set of predicate schemas,
\(\mathcal F\) is a finite set of integer-valued resource functions,
\(\mathcal A\) is a finite set of action schemas, and
\(\mathcal T\) is a type hierarchy.
Ordinary typed STRIPS is obtained when \(\mathcal F=\emptyset\).

[inline def] Let each grounded action \(a\) have Boolean preconditions
\(\operatorname{Pre}(a)\), add effects \(\operatorname{Add}(a)\), and delete
effects \(\operatorname{Del}(a)\), together with the supported bounded-integer
conditions and constant updates when numeric resources are present.

[formal def] Define a state as
\[
s=(s_{\mathcal P},s_{\mathcal F}),
\]
where \(s_{\mathcal P}\) is the set of true grounded predicate atoms and
\(s_{\mathcal F}\) assigns an integer value to each grounded resource function.

[inline def] Define a grounded action \(a\) as applicable in \(s\) when its
Boolean and numeric preconditions hold, and denote the resulting deterministic
state by
\[
\gamma_D(s,a).
\]
For the Boolean component,
\[
\gamma_D(s,a)_{\mathcal P}
=
\bigl(s_{\mathcal P}\setminus\operatorname{Del}(a)\bigr)
\cup \operatorname{Add}(a).
\]

[formal def] Define a planning instance as
\[
I=\langle D,O_I,s_I^0,G_I\rangle,
\]
where \(O_I\) is a finite set of typed objects,
\(s_I^0\) is the initial state, and
\(G_I\) is an achievement condition.

[inline def] Define a classical plan as a sequence of grounded actions
\[
\pi=\langle a_1,\ldots,a_n\rangle
\]
such that every \(a_i\) is applicable in
\[
s_{i-1},
\qquad
s_i=\gamma_D(s_{i-1},a_i),
\]
and the final state satisfies the goal,
\[
s_n\models G_I.
\]

[inline def] Define a generalized-planning task as
\[
\mathcal G=\langle D,\mathcal I\rangle,
\]
where
\[
\mathcal I=\{I_1,\ldots,I_k\}
\]
is a family of planning instances sharing the same domain \(D\).

[inline def] Define a generalized plan as a single parameterised policy or
program intended to provide reusable behaviour across the instances in
\(\mathcal I\), rather than a separate grounded action sequence for each
instance.


\subsection{Achievement and Temporally Extended Goals}

[formal def] Define a finite execution trace as
\[
\xi=s_0s_1\ldots s_n,
\]
where \(s_i\) is the state observed after the \(i\)-th environment-changing
action.

[inline def] Let \(AP\) be a finite set of state propositions, where each
\(a\in AP\) denotes either a grounded predicate condition or a supported
integer-equality condition.

[formal def] Define the supported goal language \(\Phi\) by
\[
\begin{aligned}
\ell &::= a \mid \neg a,\\
\varphi &::= \ell
          \mid (\varphi\land\varphi)
          \mid X\varphi
          \mid F\varphi
          \mid (\varphi U\varphi),
          \qquad a\in AP.
\end{aligned}
\]
Here \(X\) denotes strong next, \(F\) denotes eventually, and \(U\) denotes
strong until; negation is restricted to atomic propositions.

[formal def] Define finite-trace satisfaction for
\(0\leq i\leq n\) by
\[
\begin{aligned}
\xi,i\models a
&\iff a\text{ holds in }s_i,\\
\xi,i\models\neg a
&\iff a\text{ does not hold in }s_i,\\
\xi,i\models\varphi\land\psi
&\iff \xi,i\models\varphi
    \text{ and }\xi,i\models\psi,\\
\xi,i\models X\varphi
&\iff i<n
    \text{ and }\xi,i+1\models\varphi,\\
\xi,i\models F\varphi
&\iff \exists j\in[i,n]\;:\;\xi,j\models\varphi,\\
\xi,i\models\varphi U\psi
&\iff
\exists j\in[i,n]\;:\;
\xi,j\models\psi
\text{ and }
\forall k\in[i,j)\; \xi,k\models\varphi.
\end{aligned}
\]

[formal def] Define the achievement-goal fragment
\[
\Phi_{\mathrm{ach}}\subset\Phi
\]
by
\[
\psi ::= \ell \mid (\psi\land\psi).
\]

[inline def] Define an achievement goal as
\[
\varphi\in\Phi_{\mathrm{ach}},
\]
and say that a terminating execution
\(\xi=s_0\ldots s_n\) satisfies the achievement goal when
\[
\xi,n\models\varphi.
\]

[inline def] Define the temporally extended goal fragment as
\[
\Phi_{\mathrm{teg}}
=
\Phi\setminus\Phi_{\mathrm{ach}},
\]
and say that an execution trace \(\xi\) satisfies a TEG
\(\varphi\in\Phi_{\mathrm{teg}}\) when
\[
\xi,0\models\varphi.
\]
